\chapter{Petitions and Referenda}

\section{Petitions}

\subsection{Petition Purpose}
Petitions may be initiated by any general member of the Association
to propose referenda or to call for special meetings of the members 
to discuss specific issues related to the Association.

The goal of the petition process is to provide a mechanism for 
general members to officially raise concerns, propose changes, 
or seek collective decision-making on matters of importance 
to the Association.

\subsection{Petiton Scope}
Petitions may be used to call in a special meeting of the members 
to discuss and vote on a specific issue or to trigger a 
referendum to vote on a specific issue related to the Association.

\subsection{Petition Process}\label{petition-process}

\begin{enumerate}
    \item (Submission) A petition must be submitted in writing to the 
    Board of Executives and outline the specific issue or proposal, 
    actionable clauses, and supporting documents if applicable.

    \item (Review) Upon receipt of a petition, the Board shall review the
    petition to ensure that it meets the necessary requirements and is
    relevant to the Association's affairs and issue a response to the petitioner
    within two calendar weeks.

    \item (Signature Requirement) The petition will be then served to
     all general members of the Association and must be signed by at least
     10\% of the general members of the Association within 4 calendar
     months of the petition's serving date to be considered valid.

    \item (Verification) Upon receipt of a petition, the Board shall
    verify the validity of the petition, including confirming that it
    meets the signature requirement.

    \item (Action) If a petition is deemed valid, the Board shall take
    appropriate action based on the nature of the petition, which may
    include scheduling a special meeting or initiating a referendum.

    \item (Invalid Petitions) If a petition is deemed invalid, the Board shall
    provide a written explanation to the petitioner within ten (10) calendar
    days of the verification process, outlining the reasons for the 
    invalidation and any necessary steps for correction.
\end{enumerate}

\subsection{Non-resubmission}
A petition that did not meet the signature requirement within the specified
time frame shall not be resubmitted for the same issue for a period of 
six (6) calendar months from the date of the petition's serving date.

\section{Referenda}

\subsection{Referendum Purpose}
Referenda may be initiated by the Board of Executives or by a valid
petition to allow general members to vote on specific issues related
to the Association. The purpose of referenda is to enable direct
democratic decision-making on important matters that affect the
Association and its members.

\subsection{Referendum Scope}
Referenda may be held to vote on the following matters:
\begin{enumerate}

    \item Amendments to the constitution,

    \item Impeachment of a senior executive, and

    \item Other matters as determined by the Board of Executives 
    or through a valid petition.
\end{enumerate}
The exact details on constituional amendments and impeachment referenda 
are outlined in Articles 7, 8, and 9 of this constitution.

\subsection{Referendum Process}
\begin{enumerate}

    \item (Notice) The Board of Executives shall provide reasonable
    notice to general members about the upcoming referendum, including
    the issue being voted on, the rationale for the referendum, and 
    instructions on how to participate in the voting process.

    \item (Voting) General members shall have the opportunity to cast
    their votes in a secure and confidential manner, as determined by
    the Board of Executives. The voting period shall be open for a
    reasonable duration to allow all general members to participate.

    \item (Outcome) The outcome of a referendum shall be determined by
    a simple majority vote of the general members who participated in
    the voting process. The results shall be announced promptly after
    the conclusion of the voting period.

\end{enumerate}